\section{Control analógico y medida de rocío del tratamiento: ENV-34}
\label{sec:ambiente-env34}
\textbf{Selección y distribución A/B.} Se seleccionan dos controladores WAGO PFC200 de segunda generación \texttt{750-8212}, uno por banco, alimentados desde el riel ambiental de su banco mediante el ORING A/B seleccionado en ENV-40. El presupuesto ENV-34 que sigue constituye la base de receptores trasladados, sin duplicar nodos. Cada nodo incorpora un \texttt{750-455} de cuatro entradas 4--20 mA, dos \texttt{750-559} de cuatro salidas 0--10 V cada uno, un \texttt{750-430} de ocho entradas digitales 24 V y un módulo final \texttt{750-600}. Se reserva en cada nodo la incorporación del \texttt{750-479} diferencial de diagnóstico DC seleccionado en ENE-35; sus canales no se cuentan como entradas 4--20 mA. El orden de obra es CPU, AI de caudal/rocío, AI de corriente, AO de ventiladores, AO de válvulas, DI y final de bus. No se conectan positivos de entrada A/B por señales analógicas. El retorno secundario común necesario para ENV-40 se define en su esquema, conservando independientes las referencias de medida aisladas de CEN-41. Cada controlador regula las dos rutas de su banco; no controla por sí solo los dos bancos ni constituye redundancia de seguridad funcional \parencite[pp.~3--4]{lote34-pfc200}\parencite{lote34-end}.

La CPU admite 24 V DC $-25/+30\,\%$, publica 550 mA de entrada \emph{típicos} a carga nominal y 1.700 mA disponibles para el sistema. Sus 0--55\,$^{\circ}$C, IP20 y ausencia de condensación requieren armario interior seco de control, fuera de los bancos y del retorno de gas. Se fija para ese armario ambiente de proyecto de 24\,$^{\circ}$C; no se atribuye esa temperatura a equipos existentes. Se llevan 24 V/0 V a los bornes identificados de alimentación de sistema y de campo, con ramales protegidos, PE y referencia de señales definidos por el manual y el unifilar final. La referencia común de canales de un mismo AI/AO no equivale a aislamiento entre canales; los 500 V publicados corresponden a sistema/campo. El esquema definitivo de protecciones, sección de conductor, tierra y máxima corriente de arranque conserva su cálculo independiente.

\textbf{Mapa de señales por banco.} En el AI \texttt{750-455}, canal 1 recibe DPT de ruta 1, canal 2 DPT de ruta 2, canal 3 rocío de ruta 1 y canal 4 rocío de ruta 2. Son ocho entradas utilizadas en los dos bancos. El primer AO \texttt{750-559} dedica sus cuatro canales a los cuatro motores EC del banco, uno por motor. En el segundo AO se destinan los canales 1 y 2 a las dos RVAZ4-24A; quedan 3 y 4 de reserva, sin atribuirles una función instalada. Son doce salidas asignadas y cuatro libres en total. En cada DI se reservan 1/2 para estados auxiliares de las dos rutas de la cadena independiente, 3/4 para alarmas ML de los dos calefactores, 5/6 para estados optoacoplados LEM de ENE-35 y 7/8 para guardas auxiliares de las rutas. La interfaz aislada LEM/24 V permanece pendiente: esa reserva de canal no acredita cableado directo ni disponibilidad de una salida compatible. Toda señal de la cadena o guarda requiere contacto auxiliar independiente o interfaz seleccionada; no se deriva alimentación del contacto que ya corta una carga de seguridad \parencite[pp.~8, 11--12, 15]{lote34-ai}\parencite[pp.~7, 13]{lote34-ao}\parencite[pp.~8, 14]{lote34-di}.

Para AI y AO actuales, los bornes de señal de canales 1, 2, 3 y 4 son respectivamente 1, 5, 3 y 7; sus referencias son 2, 6, 4 y 8. Se aplican esos números a las revisiones documentadas, sin trasladar esquemas antiguos de la familia. El \texttt{750-455} tiene resistencia de entrada máxima 100 $\Omega$, resolución de catálogo de 12 bits, conversión típica de 10 ms, error máximo 0,1\,\% del valor superior de rango a 25\,$^{\circ}$C y deriva máxima 0,01\,\%/K. Se calcula conservadoramente sobre 20 mA y se incluye cuantización. El formato actual reserva B3/B2 y utiliza B1/B0 para diagnóstico; se interpreta el entero con signo y se ignoran los cuatro bits inferiores para la magnitud. No se tratan bits reservados como datos ni se aplica sin revisión el desplazamiento de manuales antiguos. El programa debe comprobar diagnóstico, rango y antigüedad antes de convertir a unidades de proceso \parencite[pp.~8, 11--12, 15]{lote34-ai}\parencite[pp.~7, 13]{lote34-ao}.

El \texttt{750-559} publica carga $\geq5$ k$\Omega$, 12 bits, conversión típica 10 ms y recuperación típica 100 ms, con error máximo 0,1\,\% a 25\,$^{\circ}$C y deriva máxima 0,01\,\%/K. Cada ventilador recibe una salida individual hacia terminal 3 YE y referencia a 2 BU, retirando el potenciómetro de fábrica conforme ENV-28. No se usa la salida 4 RD de 10 V del motor como fuente de la CPU, ni se unen dos entradas EC en una sola AO. Cada actuador recibe señal roja, 24 V negro y referencia azul del riel ambiental de su banco seleccionado en ENV-40. La impedancia de entrada EC y RVAZ4-24A no está publicada en las fuentes cotejadas: falta acreditar que cada carga real cumple $\geq5$ k$\Omega$ o seleccionar interfaz adecuada. La coincidencia 0--10 V no cierra esa compatibilidad. Antes de operar se verifica además sentido de apertura/cierre de la pareja válvula/actuador; se expresa la demanda de frío en porcentaje de apertura y no se presupone su polaridad eléctrica.

\textbf{Rocío seleccionado y circuito.} Se seleccionan cuatro Vaisala \texttt{DMT132} y cuatro bridas de conducto \texttt{DM240FA}, uno por ruta. Se montan en el aire limpio tratado después de recalentamiento/ventiladores y antes de ingresar al banco; el recalentamiento sin aporte de agua conserva el rocío. La sonda no se instala en retorno con H$_2$, dentro del serpentín mojado ni expuesta a gotas. El rango de catálogo es $-30$ a $+50\,^{\circ}$C de rocío, gas no corrosivo, temperatura $-30$ a $+50\,^{\circ}$C y salida de dos hilos 4--20 mA; IP65 no convierte la medida en inmune a condensación. La ficha publica $\pm1\,^{\circ}$C a temperatura de gas de 20\,$^{\circ}$C para rocío $-3$ a $20\,^{\circ}$C y una gráfica de exactitud en otras temperaturas. El punto previsto de aire a 23\,$^{\circ}$C y rocío 8,5\,$^{\circ}$C está dentro de la banda gráfica de $\pm1\,^{\circ}$C; no se extiende esa cifra a toda la banda del transmisor \parencite[pp.~1--2]{lote34-dmt-ficha}.

Se fija la escala de proyecto $-10$ a $+30\,^{\circ}$C mediante \texttt{ASEL TDF -10 30}: 4 mA representan $-10\,^{\circ}$C, 20 mA $30\,^{\circ}$C y 8,5\,$^{\circ}$C representan 11,4 mA. Se usa rocío a la presión real del conducto, no el rocío convertido desde una línea de aire comprimido. Se registra la presión configurada y se verifica contra la presión absoluta local en puesta en servicio. Se fija notificación \texttt{AERR 3.5}; 3,5 mA queda por debajo del rango válido del AI. Vaisala advierte que una falla grave puede impedir emitir la corriente de error: no se considera diagnóstico exhaustivo ni permiso de recarga. Configuración, lectura de retorno y prueba de corriente se guardan en el acta de cada transmisor \parencite{lote34-dmt-escala,lote34-dmt-error}.

En puerto I, pin 1 marrón recibe positivo 24 V; pin 3 azul entrega retorno de lazo hacia el AI; la referencia del AI vuelve a 0 V del mismo camino. Los pines 2/4 no se usan como salidas adicionales. La resistencia total externa, incluida entrada y cable, debe permanecer bajo 500 $\Omega$ con alimentación entre 20 y 28 V. Para la reserva de tendido propio de 10 m ida, par de cobre de sección efectiva mínima 0,25 mm$^2$ y resistividad de cálculo 0,0175 $\Omega$ mm$^2$/m, el lazo de 20 m añade 1,4 $\Omega$: $100+1,4=101,4\;\Omega$. El cable/conector M8 y su sección efectiva se especificarán en compra antes de aceptar ese presupuesto; no se asigna esa sección a un accesorio sin ficha. También se debe acreditar tensión mínima en bornes durante transferencia: si cae por debajo de 20 V, la cota OEM pasa a 100 $\Omega$ y el presupuesto anterior deja de ser admisible. Se contemplan cuatro conjuntos de conexión, sin dar por resuelta su referencia comercial \parencite[p.~2]{lote34-dmt-quick}.

\textbf{Exactitud, dinámica y aceptación.} Como cota propia del AI a 24\,$^{\circ}$C, se reserva $0,001(20)+0,0001(20)(1)=0,022$ mA, más 0,008 mA conservadores por cuantización. Esta cota de 0,030 mA equivale a 0,075\,$^{\circ}$C para escala de rocío de 40 K/16 mA. La salida DMT agrega $\pm0,05\,\%$ de escala y deriva típica $0,005\,\%/^{\circ}$C respecto a 20\,$^{\circ}$C; su resolución es 0,002 mA. Reservando 0,013 mA para esa salida a 23\,$^{\circ}$C, se obtiene presupuesto instrumental aditivo de $1+40(0,030+0,013)/16=1,1075\,^{\circ}$C, antes de instalación, presión y calibración. Es presupuesto propio con deriva \emph{típica} del DMT, no incertidumbre certificada ni garantía de fabricante. La consigna 8,5\,$^{\circ}$C medida no asegura un rocío real m\'aximo de 8,5\,$^{\circ}$C: si ese m\'aximo es el criterio de proceso, la cota instrumental exige lectura no mayor a 7,3925\,$^{\circ}$C y todavía margen de instalación. El serpentín y la fuente deben recalcularse para esa condición antes de adoptar una consigna autorizada. No se conserva una potencia calculada a 8,5\,$^{\circ}$C mientras se exige salida efectiva inferior.

El arranque publicado DMT es 3 s. Los tiempos de respuesta típicos al 90\,\% son 40 s al subir de $-14$ a $3\,^{\circ}$C y 85 s al bajar de $3$ a $-14\,^{\circ}$C, a gas 20\,$^{\circ}$C y 1 bar; no se presentan como garantía en 8,5\,$^{\circ}$C. Se propone tarea ordinaria de control cada 100 ms, con registro cada segundo y vigilancia de antigüedad del dato. Es configuración y criterio de ensayo, sin prueba de tiempo ejecutada. El caudal usa el DPT configurado a 1 s y escala propia prevista 0--1.800 m$^3$/h, compatible con una consigna de 900 m$^3$/h; la lectura nominal sería 12 mA. Esa escala debe verificarse en el DPT antes de usarla y no sustituye el coeficiente de iris ni calibración del caudal. La contribución propia máxima reservada del AI equivale a $1.800(0,030)/16=3,375$ m$^3$/h, que se añade a sensor, iris y montaje, no se confunde con el error total del sistema.

Se desarrollan dos lazos PI ordinarios separados por ruta: caudal hacia los dos EC individuales y rocío hacia la apertura de agua. Se limita cada salida a 0--10 V, se congela integración cuando una ruta está deshabilitada y se aplica antioscilación por saturación mediante seguimiento de la salida limitada. Las ganancias y límites de rampa se fijarán por ensayo del conjunto húmedo seleccionado; no se inventan constantes sin planta. Al iniciar, perder diagnóstico, recibir rango inválido o dato antiguo se descarta el integrador, se enclava alarma y se retira demanda de calor/recarga mediante su cadena independiente. No se presume que AO cero detenga todo motor, cierre la válvula o suprima gas. La cadena decide ventilación/purga segura y corte de recarga aunque CPU, programa o comunicaciones fallen; el control ordinario sólo lee auxiliares. El TTC25 conserva su regulación de recalentamiento seleccionada y sus protecciones de caudal/ML: las dos salidas AO libres no equivalen a una nueva regulación de temperatura ya instalada. El reingreso después de falla requiere estado seguro confirmado y rearme previsto de la cadena, no reinicio automático de PLC.

\textbf{Demanda DC e implantación.} Por nodo, AI de rocío/caudal 65 mA, dos AO $2(125)$ mA y DI 17 mA requieren 332 mA a 5 V; con AI diferencial ENE-35 de 100 mA son 432 mA/2,160 W, bajo los 1.700 mA de sistema. Es consumo de bus, no 432 mA a 24 V. Se reserva provisionalmente 1 A/24 W por CPU y módulos completos como cota de proyecto, incluyendo bus y pérdidas, sin convertirla en máximo OEM ni sumar de nuevo los 2,160 W. El máximo/arranque de CPU, accesorios y reserva efectiva de esa cota deben verificarse. Los dos DMT por camino reservan 30 mA cada uno para presupuesto conservador de lazo: 1,44 W/0,060 A a 24 V, aunque error elegido sea 3,5 mA y el máximo normal de escala 20 mA. Ocho DI a 2,8 mA típicos/canal agregan 0,5376 W/0,0224 A por camino. El incremento propio total es 25,9776 W/1,0824 A por camino, 51,9552 W/2,1648 A para A/B. Antes de los receptores de ENE-35, la reserva DC acumulada pasa de 101,6 a 153,5552 W, 76,7776 W por QUINT. Se integra en conversión/pérdidas de fuente y autonomía; no se suma como carga AC independiente ni se declara autonomía acreditada.

Se fija reserva adicional por nodo de armario interior de 600 por 600 por 250 mm, riel DIN útil mínimo 300 mm para el conjunto, y canaletas separadas para señal/energía. Son dimensiones de obra; falta seleccionar armario, ventilación, borneras, fusibles y accesorios de montaje reales. Se mantiene la reserva de UTA, plenum y estación exterior vigente, sin reducir recorridos de mantenimiento ni redistribuir el predio. La recepción prevista comprende inventario/revisiones de módulos, esquemas de lazos, inyección de corriente válida y fuera de rango, pérdida de energía/dato, sentido y carga real de AO, calibración de caudal/rocío, estabilidad PI y prueba conjunta de transferencia/cadena independiente. ENV-34 concreta control y medida; RT-06.03, RT-06.09 y RT-06.13 permanecen pendientes por implantación completa, compatibilidad de cargas AO, circuitos protegidos, serpentín/fuente/bomba, incertidumbre de proceso y seguridad de gas todavía abiertas.
